\begin{table}[htbp]
\centering
\footnotesize
\caption{跨病原截距比与纯报告增长零假设的比较(Holm 校正 $p$ 值)。}
\label{tab:s6}
\begin{tabular}{lrllllrrrr}
\toprule
病原,窗口 & $N$ & $q_1$ [95\% CI] & $q_2$ [95\% CI] & 零假设 $q_1$ & 零假设 $q_2$ & $q_1<1$ & $q_2>1$ & $q_1$ 越界 & $q_2$ 越界 \\
\midrule
流感,规则 & 1{,}756 & 0.68 [0.58, 0.81] & 1.61 [1.41, 1.87] & 0.63--0.89 & 1.13--1.58 & 0.020 & 0.034 & 1.000 & 1.000 \\
流感,日历 & 1{,}267 & 0.69 [0.53, 0.85] & 1.54 [1.27, 1.88] & 0.63--0.89 & 1.13--1.58 & 0.020 & 0.020 & 1.000 & 1.000 \\
流感,增长期 & 2{,}707 & 0.57 [0.47, 0.71] & 1.88 [1.55, 2.45] & 0.68--0.88 & 1.33--1.58 & 0.034 & 0.020 & 1.000 & 1.000 \\
流感,增长期(2025-26 季) & 663 & 0.64 [0.45, 0.84] & 1.56 [1.28, 2.33] & 0.66--0.89 & 1.25--1.57 & 0.116 & 0.075 & 1.000 & 1.000 \\
新冠,日历 & 1{,}479 & 0.73 [$-$0.32, 0.91] & 1.53 [0.67, 1.87] & 0.61--0.83 & 1.42--1.67 & 0.287 & 0.616 & 1.000 & 1.000 \\
新冠,增长期 & 2{,}862 & 0.78 [0.61, 1.01] & 1.41 [1.21, 2.04] & 0.60--0.84 & 1.40--1.66 & 0.616 & 0.154 & 1.000 & 1.000 \\
合胞病毒,日历 & 494 & 0.70 [0.48, 0.93] & 1.57 [1.19, 2.09] & 0.17--1.11 & 1.01--3.02 & 0.154 & 0.228 & 1.000 & 1.000 \\
合胞病毒,增长期 & 1{,}073 & 0.64 [0.11, 0.87] & 1.69 [$-$0.22, 3.06] & 0.57--1.14 & 0.89--1.69 & 0.182 & 0.660 & 1.000 & 1.000 \\
\bottomrule
\end{tabular}
\par\smallskip
\begin{minipage}{0.95\linewidth}\scriptsize $q_1=a_{\text{单周}}/a_{\text{原权重}}$, $q_2=a_{\text{三周等权}}/a_{\text{原权重}}$; 共享双向自助法, $B=2000$。零假设范围来自无传播异质性的增长期零模拟（三种世代间隔、三种入院延迟、四档报告噪声；规则窗口行沿用流感日历窗口范围）。后四列为单侧自助 $p$ 值，对全部 32 个比值检验做 Holm 校正；越界检验比值是否落在零假设范围之外。部分下置信界为负值（如新冠日历窗口 $q_1$ 及合胞病毒增长期 $q_2$），系因少数自助抽样中基准分母截距出现负值所致。\end{minipage}
\end{table}
